\section{Contoh Program}
\label{sec:contoh-bab16}

Bagian ini menyajikan tiga program yang mengerjakan pekerjaan
\emph{sama persis}, yaitu menjumlahkan sepuluh elemen larik lalu
mengalikan jumlahnya dengan delapan. Larik berisi
$10, 20, \dots, 100$, sehingga jumlahnya $550$ dan hasil akhirnya
$4400$ atau \hex{1130}. Ketiga program dibedakan hanya oleh cara
mengerjakannya, sehingga keduanya dapat dibandingkan secara adil.

\subsection{Contoh 1: Versi Naif}

Kode~\ref{lst:program-naif} adalah versi yang ditulis tanpa
memikirkan efisiensi. Program ini benar dan mudah dibaca, tetapi
memuat beberapa pemborosan.

\lstinputlisting[
    language=MASM,
    caption={Versi naif: \file{naif.asm}.},
    label={lst:program-naif}
]{\CodeDir/naif.asm}

Empat inefisiensi dapat ditemukan pada program tersebut.

\begin{enumerate}
    \item Instruksi \lstasm|ADD AX, larik[SI]| memakai pengalamatan
          langsung berindeks, sehingga alamat efektifnya memerlukan
          8 siklus, bukan 5 siklus seperti \lstasm|[SI]|.
    \item \lstasm|MOV AX, 0| berukuran tiga byte, padahal ada cara
          yang hanya dua byte.
    \item Hasil akhir ditulis ke memori (\code{jumlah}) lalu dibaca
          kembali, padahal nilainya sudah ada di \reg{AX}. Ini
          menambah dua akses memori yang tidak perlu.
    \item \lstasm|MUL BX| untuk mengalikan dengan delapan adalah
          pilihan paling mahal di seluruh program: 118--133 siklus
          untuk pekerjaan yang dapat diselesaikan dalam enam siklus.
\end{enumerate}

\subsection{Contoh 2: Versi Teroptimasi}

Kode~\ref{lst:program-teroptimasi} memperbaiki keempat hal tersebut
tanpa mengubah hasil akhir.

\lstinputlisting[
    language=MASM,
    caption={Versi teroptimasi: \file{teroptimasi.asm}.},
    label={lst:program-teroptimasi}
]{\CodeDir/teroptimasi.asm}

Empat perbaikan yang diterapkan adalah sebagai berikut.

\begin{itemize}
    \item \lstasm|MOV SI, OFFSET larik| membuat \reg{SI} menunjuk
          langsung ke data, sehingga pengalamatan
          \lstasm|ADD AX, [SI]| hanya memerlukan 5 siklus alamat
          efektif.
    \item \lstasm|XOR AX, AX| mengosongkan akumulator dalam dua byte
          dan tiga siklus.
    \item Nilai antara tidak lagi bolak-balik ke memori; \reg{AX}
          langsung dipakai sebagai akumulator sampai akhir.
    \item \lstasm|MUL BX| diganti tiga kali \lstasm|SHL AX, 1|.
          Totalnya enam siklus, yaitu sekitar dua puluh kali lebih
          cepat daripada satu instruksi \asm{MUL}.
\end{itemize}

\subsection{Contoh 3: Versi Teroptimasi dengan \eng{Loop Unrolling}}

Kode~\ref{lst:program-unroll} menambahkan satu teknik lagi:
isi perulangan digandakan sehingga satu iterasi mengerjakan dua
elemen. Jumlah iterasi turun dari sepuluh menjadi lima.

\lstinputlisting[
    language=MASM,
    caption={Versi teroptimasi dengan \eng{unrolling}:
    \file{teroptimasi-unroll.asm}.},
    label={lst:program-unroll}
]{\CodeDir/teroptimasi-unroll.asm}

\subsection{Perbandingan Hasil Disassembly}

Gambar~\ref{fig:disasm-bandingan} menampilkan potongan hasil
\eng{disassembly} kedua versi utama secara berdampingan. Keduanya
dibongkar dengan asumsi program dimuat pada alamat
\hex{1D2F:0100}, larik berada pada offset \hex{01F4}, variabel
\code{jumlah} pada offset \hex{0208}, dan variabel \code{hasil} pada
offset \hex{020A}.

\begin{figure}[H]
    \centering
    \begin{minipage}[t]{0.49\textwidth}
        \centering
        \textbf{Versi naif}\\[0.3em]
        \begin{lstlisting}[language=MASM, numbers=none, aboveskip=0pt]
1D2F:0100 B90A00   MOV  CX,000A
1D2F:0103 BE0000   MOV  SI,0000
1D2F:0106 B80000   MOV  AX,0000
1D2F:0109 0384F401 ADD  AX,[SI+01F4]
1D2F:010D 83C602   ADD  SI,+02
1D2F:0110 E2F7     LOOP 0109
1D2F:0112 89060802 MOV  [0208],AX
1D2F:0116 8B060802 MOV  AX,[0208]
1D2F:011A BB0800   MOV  BX,0008
1D2F:011D F7E3     MUL  BX
1D2F:011F 89060A02 MOV  [020A],AX
        \end{lstlisting}
        {\scriptsize 11 instruksi, 35 byte}
    \end{minipage}
    \hfill
    \begin{minipage}[t]{0.49\textwidth}
        \centering
        \textbf{Versi teroptimasi}\\[0.3em]
        \begin{lstlisting}[language=MASM, numbers=none, aboveskip=0pt]
1D2F:0100 B90A00   MOV  CX,000A
1D2F:0103 BEF401   MOV  SI,01F4
1D2F:0106 31C0     XOR  AX,AX
1D2F:0108 0304     ADD  AX,[SI]
1D2F:010A 83C602   ADD  SI,+02
1D2F:010D E2F9     LOOP 0108
1D2F:010F D1E0     SHL  AX,1
1D2F:0111 D1E0     SHL  AX,1
1D2F:0113 D1E0     SHL  AX,1
1D2F:0115 89060A02 MOV  [020A],AX
        \end{lstlisting}
        {\scriptsize 10 instruksi, 25 byte}
    \end{minipage}
    \caption{Perbandingan hasil \eng{disassembly} versi naif dan
    versi teroptimasi.}
    \label{fig:disasm-bandingan}
\end{figure}

Tiga perbedaan penting tampak pada kedua kolom tersebut.

\begin{itemize}
    \item \lstasm|ADD AX, [SI+01F4]| pada versi naif mengandung
          \eng{displacement} \hex{01F4}: nama simbolik \code{larik}
          hilang dan berubah menjadi angka. Versi teroptimasi
          menghasilkan \lstasm|ADD AX, [SI]| yang lebih pendek dua
          byte.
    \item \lstasm|MOV AX, 0000| (3 byte) digantikan
          \lstasm|XOR AX, AX| (2 byte).
    \item Urutan \lstasm|MOV [0208],AX|, \lstasm|MOV AX,[0208]|,
          \lstasm|MOV BX,0008|, \lstasm|MUL BX| digantikan oleh tiga
          \lstasm|SHL AX, 1|. Perhatikan bahwa hasil
          \eng{disassembly} sama sekali tidak memberi tahu mengapa
          perubahan itu dilakukan; penjelasannya hanya ada pada
          perbandingan metrik.
\end{itemize}

\subsection{Perbandingan Hasil Pengukuran}

Tabel~\ref{tab:hasil-pengukuran-16} merangkum perbandingan ketiga
versi. Angka siklus dihitung memakai perkiraan pada
Tabel~\ref{tab:siklus-8086-perkiraan} dengan asumsi tanpa
\eng{wait state} dan \eng{prefetch queue} selalu terisi. Bagian yang
diukur adalah bagian yang ditandai komentar pada ketiga berkas.

\begin{table}[H]
    \centering
    \caption{Perbandingan hasil pengukuran versi naif dan versi
    teroptimasi.}
    \label{tab:hasil-pengukuran-16}
    \small
    \begin{tabular}{@{}p{3.6cm}cccc@{}}
        \toprule
        \textbf{Versi} & \textbf{Instruksi} & \textbf{Ukuran} &
        \textbf{Perkiraan} & \textbf{Akses} \\
        & \textbf{dieksekusi} & \textbf{kode} &
        \textbf{siklus} & \textbf{memori} \\
        \midrule
        Naif & 38 & 35 byte & $\approx$546 & 13 \\
        \addlinespace
        Teroptimasi & 37 & 25 byte & $\approx$370 & 11 \\
        \addlinespace
        Teroptimasi + \eng{unrolling} & 27 & 28 byte & $\approx$280
        & 11 \\
        \bottomrule
    \end{tabular}
\end{table}

Tabel~\ref{tab:hasil-pengukuran-16} menyimpan pelajaran penting.
Pertama, jumlah instruksi yang dieksekusi hampir tidak berubah
antara versi naif dan versi teroptimasi (38 berbanding 37), padahal
perkiraan siklusnya turun sekitar 32 persen. Artinya, \important{jumlah
instruksi saja adalah metrik yang menyesatkan}; yang menentukan
adalah \emph{jenis} instruksi dan banyaknya akses memori. Kedua,
versi \eng{unrolling} paling cepat tetapi ukurannya bertambah tiga
byte dibandingkan versi teroptimasi; inilah trade-off yang dibahas
pada Subbab~\ref{sec:optimasi-instruksi}. Ketiga, angka-angka ini
baru perkiraan. Pembuktiannya tetap memerlukan pengukuran waktu pada
mesin, seperti yang akan dikerjakan pada praktikum.
